\section{Limitations}
\label{sec:limitations}

This article restricts the implementation and verification to immutable,
finite lists of integers represented using the
\texttt{stainless.collection.List} data type. The focus is on
\textbf{correctness}, not on performance or scalability. Our summation
and accumulation models follow a \textbf{recursive definition}, aligned
with mathematical formalism. However, this approach may introduce
performance limitations in practical applications involving large lists.

\subsection{Overflow and Memory Limits Are Out of Scope}

By using \texttt{BigInt} and immutable lists, the model assumes
unbounded integer arithmetic and infinite list capacity. This choice
avoids overflow and out-of-memory errors, but it does not reflect the
constraints of fixed-size integer types or limited system memory in
real-world environments.

\subsection{Side Effects Are Excluded}

All list operations are pure and referentially transparent. Mutation,
I/O, and performance overhead are outside the scope of this model.

\subsection{No Parallelism or Laziness}

Unlike streaming libraries or lazy sequences, this model is strictly
eager and sequential, without support for parallel computation or lazy
evaluation.

\subsection{Limitations Imposed by Stainless Verification Tool}

Due to current limitations in the Scala Stainless verifier (version
0.9.8.8), formal proofs must often rely on concrete numeric types such as
\texttt{BigInt}. Stainless does not yet fully support generic numeric
abstractions or type classes like \texttt{Numeric[T]}, which hinders
verification of implementations parameterized over arbitrary numeric
types.

As a result, while the mathematical properties in this work conceptually
apply to any numeric domain satisfying the required algebraic laws,
practical verification is constrained to \texttt{BigInt}. Overcoming
these tool limitations is an important direction for future
enhancements, enabling broader generality and more flexible formal
verification.

\subsection{Scope of Correctness}

This article emphasizes the \textbf{mathematical correctness} of
recursive definitions and verified properties, rather than runtime
behavior or system-level efficiency.

The use of \texttt{BigInt} and conceptually unbounded lists abstracts
away concerns like stack overflows, memory usage, and execution time. It
also circumvents limitations of the current version of Scala Stainless
with respect to generic numeric reasoning.

While limiting practical use in some contexts, these assumptions
maintain the focus on proving functional correctness as defined by
recursive specifications.

Future work may include developing alternative implementations of these
data structures that explicitly address real-world constraints, such as
bounded memory and side effects, alongside formal proofs establishing
their equivalence with the current, mathematically rigorous model.
